\chapter{Models, Explanation, Prediction, and Laws}
\label{ch:models}

\begin{chapterguide}
Models are treated as selective mediators between theory and world. The
chapter explains how idealized or literally false models can be scientifically
useful, compares major accounts of explanation, and asks what laws of nature
are doing in actual scientific practice. It then connects reduction and
emergence to the idea of scale-specific constraints.
\end{chapterguide}

Scientific theories are rarely used as complete descriptions. Scientists build
models: simplified systems, equations, diagrams, simulations, analogies,
mechanistic sketches, statistical representations, and idealized populations.
A model may contain frictionless planes, perfectly rational agents, isolated
genes, representative households, point particles, or infinite populations.
None of these may exist exactly. Yet the models can reveal dependencies,
generate predictions, and guide interventions.

The philosophical mistake is to treat usefulness and literal truth as
competitors. The right question is what a model represents, under which
transformations, and for which purpose.

\section{Models as mediators}

Mary Morgan and Margaret Morrison describe models as mediators that are
neither simply derived from theory nor copied from the world
\citep{morganmorrison1999}. A model can borrow resources from several
theories, data sources, and experimental practices. It provides a manipulable
object for reasoning.

Ronald Giere's account is similarly selective. A model represents a real system
to a degree and for a purpose when the relevant similarities are identified
\citep{giere1988,giere2004}. A model of a population may represent its age
distribution but not its cultural history. A model of a pendulum may represent
period under small-angle conditions while ignoring air resistance. The model's
success is conditional on what is preserved.

A representation relation is therefore not binary. We can write informally:
\[
  M \xrightarrow[\text{purpose, scale, conditions}]{\text{represents}}
  T,
\]
where \(M\) is a model and \(T\) is a target system. The arrow is labelled
because the same model can represent different features of a target under
different purposes. The representation is strong when relevant relations are
preserved and weak when the model changes the target's behaviour in the
dimension that matters.

\section{Idealization and abstraction}

An idealization deliberately removes factors. Abstraction selects a feature
while leaving other features unspecified. These are not the same. A
frictionless plane removes friction by setting a parameter to zero. A map
abstracts from buildings when it preserves roads and elevation. An agent-based
model may abstract from individual biographies while idealizing agents as
following a simple rule.

A useful idealization should satisfy three conditions.

\begin{enumerate}
\item \textbf{Relevance:} the removed feature should not be central to the
question under investigation.
\item \textbf{Stability:} the result should remain approximately valid when the
removed feature is reintroduced within the intended range.
\item \textbf{Boundary:} the model should state where the idealization breaks
down.
\end{enumerate}

A model can be literally false and structurally accurate. For a simple example,
consider a population model with exponential growth:
\[
  N(t) = N_0 e^{rt}.
\]
Here \(N(t)\) is the population at time \(t\), \(N_0\) is the initial
population, and \(r\) is a constant net growth rate. No real population grows
with an exactly constant \(r\) forever. The model can still represent a local
growth regime, estimate a rate, or provide a baseline against which more
complex models are compared.

The danger is not idealization itself. The danger is forgetting that a result
about the model is being treated as a result about the target without a
transport argument.

\section{Prediction and explanation}

Prediction and explanation overlap but are not identical. A model can predict
an outcome using a reliable correlation without revealing why the outcome
occurs. A mechanism can explain a process without providing a precise long-run
forecast when initial conditions are unknown or the system is chaotic.

The deductive-nomological model associated with Hempel and Oppenheim treats an
explanation as a deduction of the explanandum from laws and initial conditions
\citep{hempeloppenheim1948,hempel1965}. Its great virtue is clarity: it asks
which premises entail the result. Its limitation is that logical derivation
alone does not guarantee explanatory relevance. An irrelevant law can be
included in a valid deduction, and probabilistic explanations do not entail
their outcomes.

Unificationist approaches evaluate explanations by how many phenomena they
derive from a small set of argument patterns. This captures a real value:
compressing unrelated facts into one framework can reveal a shared structure.
But simplicity can mislead when the world is heterogeneous.

Causal and mechanistic approaches focus on how a phenomenon is produced.
Salmon emphasized causal structure and statistical relevance
\citep{salmon1984}. Mechanistic philosophers analyze organized parts and
activities whose interaction produces a behavior
\citep{bechtel2005,craver2007}. These accounts explain why an intervention or
process makes a difference, not only why an equation fits.

\begin{principlebox}
\textbf{Explanatory adequacy principle.} A scientific explanation is adequate
only relative to a question. For a how-question, a mechanism may be required.
For a prediction question, a calibrated statistical model may be enough. For a
why-question about a general pattern, unification or a law-like relation may
matter. No single form of explanation is the universal measure of science.
\end{principlebox}

\section{Laws of nature and local regularities}

A law is often imagined as a universal sentence that describes what nature does
in every circumstance. Actual science is more complicated. Laws may be exact
ideal relations, approximate regularities, symmetry principles, constraints on
possibilities, or statements that hold only under a ceteris paribus clause.

Nancy Cartwright argues that fundamental laws can be highly idealized and that
scientific explanation often proceeds through models that construct a
\enquote{simulacrum} of a phenomenon rather than by direct deduction from
universal laws \citep{cartwright1983,cartwright1999}. This is not a denial of
reality. It is a warning that a law's role in a model may differ from a law's
metaphysical status.

D. M. Armstrong's account treats laws as relations among universals
\citep{armstrong1983}. Other accounts treat them as regularities, dispositions,
or modal constraints. \RCR\ does not settle the ultimate metaphysics by
definition. It says that a law-like claim earns scientific authority when it
captures a stable constraint across a specified range of interventions and
conditions, and when deviations can be explained by additional factors rather
than ignored.

\section{Reduction and emergence}

Reduction has at least three meanings:

\begin{enumerate}
\item \textbf{Ontological reduction:} higher-level entities are nothing over
and above lower-level entities.
\item \textbf{Explanatory reduction:} higher-level laws or explanations can be
derived from lower-level laws.
\item \textbf{Methodological reduction:} a problem is best investigated by
isolating lower-level components.
\end{enumerate}

These meanings can come apart. A social institution may be constituted by
people and rules while still requiring concepts that are not useful in particle
physics. A biological pathway may be physically realized while explanations at
the molecular, cellular, and organismic levels remain distinct. Nagel's
classical model of reduction emphasizes derivation and bridge principles
\citep{nagel1961}; later work has shown that reduction in practice can be
local, approximate, and model-dependent.

Emergence is not a magic word for ignorance. A property is emergent in a
scientifically meaningful sense when it depends on lower-level organization but
has stable patterns, explanatory roles, or causal powers at a higher level.
Phase transitions provide cases where limiting behavior and scale can make
higher-level regularities more informative than a microscopic description
alone \citep{batterman2002}.

\RCR\ treats levels as nested constraint regimes. A higher-level model is not
automatically anti-realist because it cannot be replaced by a lower-level
description in practice. Nor is every higher-level regularity fundamental.
The question is whether the higher-level constraint is stable, manipulable,
and explanatory within its domain.

\section{Simulation and computational models}

Computer simulations are models executed as algorithms. They can explore
systems for which analytic solutions are unavailable, but their results are
not automatically empirical evidence. A simulation inherits assumptions from
its equations, discretization, parameter values, boundary conditions, and
code. It gains epistemic strength when these components are validated against
analytical results, historical data, experiments, or independent simulations.

A simulation can also discover emergent behavior. The discovery is not that
the code produced a surprise in a metaphysical vacuum. It is that a specified
set of rules generates a pattern that was not obvious from inspection and that
can be compared with a target system. The comparison requires a bridge between
the computational state and the phenomenon.

\section{The model's honesty condition}

Every scientific model should make four things visible:

\begin{enumerate}
\item what it treats as the target;
\item what it idealizes or leaves out;
\item which output is a prediction, explanation, or exploratory possibility;
\item how the output was validated and where it should not be transported.
\end{enumerate}

This is the \emph{model's honesty condition}. A model need not contain every
detail to be honest. It needs to distinguish internal consequences from
claims about the world.

\section*{Chapter summary}

Models mediate between theory and world. They are selective, purpose-relative,
and often idealized, but that does not make them arbitrary. A model is
scientifically valuable when it preserves relevant structure, remains stable
under appropriate changes, and states its boundary conditions. Prediction,
explanation, and intervention can come apart. Laws operate through idealized
models and local regimes. Reduction and emergence are best treated as
relations among levels of constraint rather than as all-or-nothing doctrines.

\begin{discussion}
\begin{enumerate}
\item When does an idealization become a distortion rather than a useful
abstraction?
\item Can a model explain a phenomenon if it cannot predict it?
\item Does a higher-level explanation remain objective when a lower-level
description is physically complete?
\end{enumerate}
\end{discussion}

\begin{furtherreading}
See Morgan and Morrison (1999), Giere (1988, 2004), Cartwright (1983, 1999),
Frigg and Hartmann (2006), Salmon (1984), and Craver (2007).
\end{furtherreading}

